% ============================================================
%  preambulo.tex  --  Matemática Financeira (Graduação Economia)
%  Referência: Bueno, Rangel, Santos -- Matemática Financeira Moderna
% ============================================================

\usepackage[T1]{fontenc}
\usepackage[utf8]{inputenc}
\usepackage{babel}
\babelprovide[main, import]{portuguese}
\usepackage{lmodern}
\usepackage{amsmath, amssymb, amsthm}
\usepackage{mathtools}

% ─── Gráficos ──────────────────────────────────────────────────────────────────
\usepackage{tikz}
\usetikzlibrary{arrows.meta, positioning, calc, patterns,
                decorations.pathreplacing, intersections}
\usepackage{pgfplots}
\pgfplotsset{compat=1.18}

% ─── Caixas e ambientes ────────────────────────────────────────────────────────
\usepackage[most, breakable]{tcolorbox}
\tcbuselibrary{theorems, skins, breakable}

% ─── Tabelas ──────────────────────────────────────────────────────────────────
\usepackage{booktabs}
\usepackage{array}
\usepackage{longtable}

% ─── Layout ───────────────────────────────────────────────────────────────────
\usepackage[margin=2.5cm, top=3cm, bottom=3cm]{geometry}
\usepackage{hyperref}
\hypersetup{colorlinks=true, linkcolor=blue!60!black,
            urlcolor=blue!70!black, citecolor=green!50!black}
\usepackage{xcolor}
\usepackage{enumitem}
\usepackage{microtype}

% ─── Cores ────────────────────────────────────────────────────────────────────
\definecolor{azul}     {RGB}{  0,  70, 127}
\definecolor{verde}    {RGB}{  0, 100,  60}
\definecolor{laranja}  {RGB}{180,  80,   0}
\definecolor{vermelho} {RGB}{160,  20,  20}
\definecolor{cinzaclaro}{RGB}{248, 248, 248}
\definecolor{azulclaro}{RGB}{220, 235, 252}

% ─── Ambientes tcolorbox ──────────────────────────────────────────────────────

% Exemplo numerado por capítulo
\newtcolorbox[auto counter, number within=chapter]{exemplo}[2][]{%
  enhanced, breakable,
  title={\textbf{Exemplo~\thetcbcounter}\quad\textmd{#2}},
  colback      = cinzaclaro,
  colframe     = azul,
  colbacktitle = azul,
  coltitle     = white,
  fonttitle    = \small,
  attach boxed title to top left={yshift=-2mm, xshift=3mm},
  boxed title style={sharp corners, colback=azul},
  #1
}

% Definição
\newtcolorbox{definicao}[1][]{%
  enhanced,
  title        = {Definição},
  colback      = azulclaro,
  colframe     = azul!80,
  colbacktitle = azul!80,
  coltitle     = white,
  fonttitle    = \bfseries\small,
  attach boxed title to top left={yshift=-2mm, xshift=3mm},
  boxed title style={sharp corners},
  #1
}

% Resultado / fórmula-chave
\newtcolorbox{resultado}[1][]{%
  enhanced,
  colback  = green!6,
  colframe = verde!80,
  boxrule  = 1.2pt,
  #1
}

% Nota / observação
\newtcolorbox{nota}[1][]{%
  enhanced, breakable,
  title        = {Nota},
  colback      = orange!6,
  colframe     = laranja!70,
  colbacktitle = laranja!70,
  coltitle     = white,
  fonttitle    = \bfseries\small,
  attach boxed title to top left={yshift=-2mm, xshift=3mm},
  boxed title style={sharp corners},
  #1
}

% Destaque moderno
\newtcolorbox{moderno}[1][]{%
  enhanced, breakable,
  title        = {Exemplo moderno},
  colback      = green!5,
  colframe     = verde!70,
  colbacktitle = verde!70,
  coltitle     = white,
  fonttitle    = \bfseries\small,
  attach boxed title to top left={yshift=-2mm, xshift=3mm},
  boxed title style={sharp corners},
  #1
}

% ─── Teoremas ─────────────────────────────────────────────────────────────────
\theoremstyle{definition}
\newtheorem{obs}{Observação}[chapter]
\newtheorem{prop}{Proposição}[chapter]

% ─── Comandos matemáticos ─────────────────────────────────────────────────────
\newcommand{\VP}{\mathrm{VP}}
\newcommand{\VF}{\mathrm{VF}}
\newcommand{\TIR}{\mathrm{TIR}}
\newcommand{\VPL}{\mathrm{VPL}}
\newcommand{\TIRM}{\mathrm{TIRM}}
\newcommand{\FAC}{\mathrm{FAC}}
\newcommand{\YTM}{\mathrm{YTM}}

% Linha do tempo genérica (uso interno nos caps)
% \fluxtimeline{n}{payment}{top_label}{top_val}
% (cada capítulo define seu próprio tikzpicture para flexibilidade)

% ─── Estilo TikZ para linhas do tempo ─────────────────────────────────────────
\tikzset{
  timeline axis/.style={->, >=Stealth, line width=0.6pt},
  tick/.style={line width=0.5pt},
  arrow down/.style={->, >=Stealth, line width=1.4pt, color=azul},
  arrow up/.style={->, >=Stealth, line width=1.4pt, color=vermelho},
  cash label/.style={font=\small},
}
